\section{Introduction}

Non-convulsive seizures occur in about 30\% of critically ill patients who undergo electroencephalography (EEG) monitoring~\cite{hirsch2021american} and are associated with  prolonged intensive care unit (ICU) stays, high healthcare costs, increased mortality, and long-term neurological impairment~\cite{rodriguez2017association,vespa2005continuous,westover2015probability}. Despite increased detection owing to widespread use of continuous EEG monitoring, optimal treatment strategies remain uncertain: Aggressive anti-seizure therapy can suppress epileptiform activity, but carries risks of sedation, hemodynamic instability, and drug toxicity, while untreated seizures may damage the brain~\cite{lalgudi2019electrographic, demarchis2016seizure, payne2014seizure, zafar2018effect, rossetti2005refractory, muhlhofer2019duration, marchi2015status, an2018variability, sutter2016status, abend2015critical}. 

This uncertainty leads to practice variation across clinicians and institutions, illustrated in Figure~\ref{fig:control-problem}. Some centers and some clinicians routinely use aggressive anti-seizure therapy, while others adopt more conservative approaches~\cite{chong2005which, rubinos2018ictal, osman2018ictal, johnson2017population, tao2020ictal, cormier2017ictal}. This variability, while reflecting uncertainty about optimal care, nevertheless provides an opportunity to learn, from observational data, which treatment strategies lead to the best outcomes. However, extracting valid causal insights from this variation is far from straightforward. Treatment decisions and disease progression influence each other over time in a feedback loop, creating two particularly acute statistical challenges. First, confounding by indication, where sicker patients are more likely to receive aggressive treatment, making treatment appear harmful when it may in fact be beneficial. Second, time-varying confounding with treatment–confounder feedback, where treatment at one time point affects subsequent disease severity, which in turn influences later treatment decisions and treatment response (e.g., electrographic improvement or adverse effects).

Experimental studies such as randomized controlled trials (RCTs) could provide valid treatment effect estimation. However, operationalizing them in critical-care seizure management is challenging due to ethical, financial, and logistical constraints. When there is substantial uncertainty about the risks related to a given treatment strategy, an RCT risks enrolling large numbers of patients into a protocol that may be ineffective or even harmful.

The two RCTs that have studied aggressive anti-seizure treatment in critically ill adults---Husain et al.~\cite{husain2018randomized} in non-convulsive status epilepticus and Ruijter et al.~\cite{ruijter2022treating} in rhythmic or periodic EEG patterns following post-anoxic encephalopathy---both reported null results for the primary outcomes. We do not claim that these null results can be explained away by a single cause. They may reflect a genuine absence of a population-average effect in the populations studied, features of the specific dosing protocols and trial designs, heterogeneity of the underlying disease, or some combination. The post-anoxic encephalopathy population in particular has a pathophysiology in which widespread ischemic injury may render any seizure-directed therapy futile regardless of protocol.

What CICADAS can contribute is orthogonal to re-litigating those null results: given the uncertainty about when and in whom EA-directed treatment is beneficial, we use causal inference on observational data---where treatment intensity varies substantially across patients, clinicians, and centers---to identify the patient subgroups and treatment intensities in which a benefit is most plausible. These subgroup-specific hypotheses are intended as inputs to the design of focused future trials rather than as contradictions of existing ones~\cite{parikh2024many}.

Because randomized controlled trials are difficult to conduct and may not capture the complexity of ICU seizure management, an alternative approach is to learn from the wealth of observational data collected in routine ICU care. However, drawing valid causal conclusions from such data is methodologically challenging. Standard regression or Kaplan–Meier analyses can be biased when treatment and disease severity influence each other over time. Modern causal inference methods—particularly the parametric g-formula—offer a principled way to overcome these challenges by modeling the joint distribution of time-varying variables and simulating counterfactual outcomes under hypothetical treatment strategies~\cite{hernan2016using, hernan2020causal, parikh2024many}.

Our work builds on two prior studies of seizure treatment in the ICU~\cite{parikh2023effects, parikh2024safe}. The first~\cite{parikh2023effects} used PK/PD modeling and matching in retrospective ICU EEG data to estimate a causal harmful effect of epileptiform activity (EA) on outcomes after accounting for EA--treatment feedback. That evidence is consistent with, but does not definitively establish, a causal role for EA: it rests on a single retrospective cohort and remains subject to residual unmeasured confounding, and the generalizability of the effect across patient subgroups and seizure etiologies is not yet known. Treating EA reduction as a provisional therapeutic target---rather than an established one---is one of the motivations for the simulation-based hypothesis-generating framework presented here. The second study~\cite{parikh2024safe} examined whether outcomes could be improved by selecting, for a given patient, among treatment patterns already observed in routine practice those associated with the highest probability of favorable outcomes. That work fit a simplified policy template to observed clinical care and introduced useful new methodology; a key limitation was that the fitted policies only modestly captured actual clinician behavior. The present work (CICADAS) advances this trajectory by moving beyond simplified policy templates to closed-loop seizure treatment strategies that more closely reflect physiologic dynamics. By integrating the parametric g-formula with mechanistic PK/PD models, disease dynamics, and simulated randomized trials, CICADAS is designed to generate physiologically plausible dynamic control targets as hypotheses for prospective evaluation rather than as direct clinical recommendations.

Here we build on this literature to present a computational framework for \emph{posing} the question of optimal seizure management as a formal causal estimation problem, and we evaluate, in simulation, whether that problem can be solved when the data-generating process is known. Our paper makes four contributions:
\begin{enumerate}
\item \textbf{Clinical motivation}: We formalize the question of optimal anti-seizure treatment in the ICU as a dynamic-treatment-regime problem, in a setting where randomized trials are impractical and the two that have been conducted were inconclusive. The framework is intended to generate hypotheses for such trials, not to guide individual treatment decisions.

\item \textbf{Methodological bridge}: We present a comprehensive tutorial with detailed mathematical derivations that connect modern causal inference theory to the practical challenges of estimating treatment effects in ICU seizure management, where treatment decisions and disease dynamics interact over time.

\item \textbf{Simulation framework with known ground truth}: We develop CICADAS (Causal Inference for Critical Care with Adherence, Disease dynamics, and Survival), a realistic simulation platform that captures seizure disease dynamics and pharmacokinetic--pharmacodynamic effects of anti-seizure medications. All analyses in this paper are performed on simulated data in which the true causal effects are known, enabling rigorous validation of the methods under controlled conditions.

\item \textbf{Exploration of the policy space}: Using the simulated cohorts, we illustrate how the framework can be used to compare a family of closed-loop policies indexed by the target threshold $\theta$, and how the comparison varies with simulated patient characteristics. These are demonstrations of what the machinery computes, within a data-generating process we specified; they are not findings about patients. External validation in real ICU EEG cohorts is required before any such comparison carries clinical meaning, and is the subject of separate, ongoing work.
\end{enumerate}

By bridging causal inference theory with clinically realistic modeling, CICADAS is intended as a framework for \textbf{targeted hypothesis generation} from observational ICU datasets. Rather than replacing randomized trials, the approach is designed to help allocate trial resources by identifying the patient subgroups and treatment intensities in which a benefit is most plausible. Whether it can do so in real data is the subject of Section~\ref{sec:external-validation} and remains untested.

The paper proceeds as follows. Section~\ref{sec:clinical-problem} states the clinical control problem
and fixes notation; Sections~\ref{sec:prelim}--\ref{sec:estimation} develop the causal framework, the
simulation model, and the estimation procedure; Section~\ref{sec:results} reports the simulation study,
including internal validation, methodological benchmarking, and robustness analyses; and
Sections~\ref{sec:discussion}--\ref{sec:conclusion} discuss what the framework can and cannot support.
Derivations, full model parameterizations, and estimation details are given in the Supplementary Material.
